\section{Problem Formulation}
\label{sec:problem}

\subsection{Notation and the index-selection setting}
Let \(q\) denote a SQL query and \(I\) a set of secondary indexes, where each
element of \(I\) is a set of columns. Query \(q\) has a set of
\emph{predicate columns} \(\pi(q)\) (the columns referenced in its
\textsc{where}/\textsc{on} clauses), extracted by a SQL parser. For a
(column, index) coverage relation we write \(\pi(q) \sqsubseteq I\) if some index
in \(I\) contains all the predicate columns of \(q\).

We model the operating cost of serving \(q\) under index configuration \(I\) as
\begin{equation}
  C(q, I) = L(q, I) + \rho\,|I|,
  \label{eq:cost}
\end{equation}
where \(L(q,I)\) is the query latency and \(\rho\,|I|\) is an amortized
maintenance/space cost charged per index per unit time. In our benchmark
\(L(q,I)\) takes one of two values depending on whether the predicate is
\emph{covered} \emph{and} \emph{selective}; in a real deployment \(L\) is
measured directly. Let \(\operatorname{sel}(c, p)\) denote the selectivity of
column \(c\) in phase \(p\) (a high-selectivity predicate benefits from an index;
a low-selectivity one does not). Then
\begin{equation}
  L(q,I) =
  \begin{cases}
    \ell_{\text{low}} & \text{if } \pi(q)\sqsubseteq I \text{ and } \pi(q)
      \text{ is selective for the phase,}\\
    \ell_{\text{high}} & \text{otherwise,}
  \end{cases}
  \label{eq:latency}
\end{equation}
with \(\ell_{\text{low}} \ll \ell_{\text{high}}\). Creating an index incurs a
one-time build cost \(\kappa\).

\subsection{The index-selection MDP}
We formulate a per-query decision problem as a Markov decision process (MDP)
\(M = (\State, \Action, P, R, \gamma)\).

\paragraph{State.}
At step \(t\) the state concatenates four components,
\begin{equation}
  s_t = \big[\,\underbrace{\phi(q_t)}_{\in\R^{74}} \;\|\;
   \underbrace{e(q_t)}_{\in\R^{384}} \;\|\;
   \underbrace{c_t}_{\in\R^{3}} \;\|\;
   \underbrace{\iota(I_t)}_{\in\R^{3}}\,\big] \in \R^{464},
  \label{eq:state}
\end{equation}
where \(\phi(q)\) is a structural feature vector (join count, presence of
\textsc{group by}/\textsc{order by}, estimated cardinality, and a hashed
bag-of-words over tables and predicate columns), \(e(q)\) is a semantic
embedding of the query text, \(c_t\) are normalized workload-context statistics
(phase index, progress, index-utilization), and \(\iota(I_t)\) is the
\emph{index-configuration state}
\begin{equation}
  \iota(I) = \big[\,|I|/B,\;\;
    \mathbb{1}[\pi(q)\sqsubseteq I],\;\;
    \max(0,(B-|I|)/B)\,\big],
  \label{eq:indexstate}
\end{equation}
with index budget \(B\). The inclusion of \(\iota\) is what makes the problem a
\emph{control} problem over index configurations rather than a one-shot
classification: the same query has different optimal actions depending on which
indexes are already present.

\paragraph{Action.}
\(\Action = \{0,1,2\}\) corresponding to
\textsc{keep}, \textsc{create\_index} (on \(\pi(q_t)\)), and
\textsc{drop\_index} (on \(\pi(q_t)\)). A structural mask
\(\mu(s_t)\in\{0,1\}^{|\Action|}\) disables infeasible actions:
\(\textsc{create}\) is allowed only if \(\pi(q_t)\neq\emptyset\) and
\(|I_t|<B\); \(\textsc{drop}\) is allowed only if an index equal to
\(\pi(q_t)\) is active.

\paragraph{Transition.}
Given \(a_t\), the index configuration updates deterministically,
\begin{equation}
  I_{t+1} =
  \begin{cases}
    I_t \cup \{\pi(q_t)\} & a_t=\textsc{create},\\
    I_t \setminus \{\pi(q_t)\} & a_t=\textsc{drop},\\
    I_t & a_t=\textsc{keep},
  \end{cases}
  \label{eq:transition}
\end{equation}
and the next query \(q_{t+1}\) is drawn from the current phase distribution.

\paragraph{Reward.}
The reward is the negative operating cost after the action, minus the build cost,
\begin{equation}
  R(s_t,a_t) = -\Big(L(q_t, I_{t+1}) + \rho|I_{t+1}|\Big)
              - \kappa\,\mathbb{1}[a_t=\textsc{create}].
  \label{eq:reward}
\end{equation}
Equation~\eqref{eq:reward} is dense, bounded, and directly tied to the quantity
an operator cares about; it makes \textsc{create} a temporally extended decision
(pay now, benefit later) and therefore justifies a discount factor \(\gamma<1\)
and a value function.

\paragraph{Objective.}
The agent seeks a policy \(\pi_\theta:\State\to\Delta(\Action)\) maximizing the
discounted return
\begin{equation}
  J(\theta) = \E_{\tau\sim\pi_\theta}\Big[\textstyle\sum_{t\ge0}\gamma^t R(s_t,a_t)\Big].
  \label{eq:objective}
\end{equation}
The optimal action-value function satisfies the Bellman optimality equation
\(Q^\star(s,a) = \E_{s'}[R(s,a) + \gamma\max_{a'} Q^\star(s',a')]\), and
\(\pi^\star(s) = \arg\max_a Q^\star(s,a)\).

\subsection{The continual setting}
Unlike standard RL, the query distribution is \emph{non-stationary}. We model the
stream as a sequence of \(T\) \emph{workload phases} (tasks)
\(\mathcal{T}_1,\dots,\mathcal{T}_T\), where phase \(\mathcal{T}_i\) induces a
query distribution \(P_i\) (and, through selectivity, a possibly different
optimal policy). The agent observes one phase at a time and may not revisit old
data unless it stores it. The \emph{continual objective} is to maximize the
average performance over all phases after the stream has been processed, while
not sacrificing adaptation speed:
\begin{equation}
  \max_\theta \; \frac{1}{T}\sum_{j=1}^{T} \E_{q\sim P_j, I}\big[\,R_\theta(q,I)\,\big].
\end{equation}

\subsection{Evaluation metrics}
Following \citet{lopezpaz2017gem}, let \(R_{i,j}\) be the agent's mean reward on
phase \(j\) after it has finished training on phase \(i\) (so \(\R\) is a
\(T\times T\) lower-triangular matrix). Let \(b_j\) be the reward of a reference
(here, a uniformly random policy) on phase \(j\). We report
\begin{align}
  \mathrm{ACC} &= \frac{1}{T}\sum_{j=1}^{T} R_{T,j}, \label{eq:acc}\\
  \mathrm{BWT} &= \frac{1}{T}\sum_{j=1}^{T} \big(R_{T,j} - R_{j,j}\big), \label{eq:bwt}\\
  \mathrm{FWT} &= \frac{1}{T}\sum_{j=1}^{T} \big(R_{j,j} - b_j\big), \label{eq:fwt}\\
  \Fforget     &= \frac{1}{T}\sum_{j=1}^{T} \Big(\max_{i\in\{j,\dots,T\}} R_{i,j} - R_{T,j}\Big). \label{eq:forgetting}
\end{align}
\(\mathrm{ACC}\) measures final average quality, \(\mathrm{BWT}<0\) indicates
forgetting of old phases, \(\mathrm{FWT}>0\) indicates positive transfer, and
\(\Fforget\ge0\) is the average drop from each phase's peak to its final
performance. A \emph{continual} agent should achieve high \(\mathrm{ACC}\),
\(\mathrm{BWT}\approx0\) and \(\Fforget\approx0\).
